\documentclass[10pt,twocolumn,letterpaper]{article}

\usepackage{cvpr}
\usepackage{times}
\usepackage{epsfig}
\usepackage{graphicx}
\usepackage{amsmath}
\usepackage{amssymb}
\usepackage{booktabs}
\usepackage{xcolor}

\usepackage[breaklinks=true,bookmarks=false]{hyperref}

\cvprfinalcopy

\def\cvprPaperID{****}
\def\httilde{\mbox{\tt\raisebox{-.5ex}{\symbol{126}}}}
\newcommand{\best}[1]{\textcolor{red!75!black}{\textbf{#1}}}
\newcommand{\second}[1]{\textcolor{blue!75!black}{\textbf{#1}}}

\setcounter{page}{1}
\begin{document}

%%%%%%%%% TITLE
\title{MultiSticker: Intent-Guided Multimodal Sticker Retrieval\\with Dialogue Memory}

\author{Anonymous Author\\
Department of Computer Science, Rice University\\
{\tt\small rl182@rice.edu}
}

\maketitle

%%%%%%%%% ABSTRACT
\begin{abstract}
Stickers are widely used in messaging platforms to express affect, humor and
intent. Retrieving the right sticker for a dialogue turn is challenging because
(i)~the candidate bank is large and dominated by animated formats, (ii)~the
correct reply depends on long-range conversational state rather than just the
last turn, and (iii)~each context typically has only one observed positive even
though many semantically equivalent stickers exist. We present
\textbf{MultiSticker}, an intent-guided sticker retrieval system that combines
three text streams---recent dialogue context, retrieved long-term session
memory, and an LLM-generated reply intent---against a frozen-or-adapted
OpenCLIP image bank covering \texttt{.png}, \texttt{.gif} and \texttt{.webm}
stickers. We compare a frozen-CLIP zero-shot baseline against a learned
retrieval head and three parameter-efficient fine-tuning variants
(text-side LoRA, image-side LoRA, dual LoRA) on the U-Sticker corpus
(8{,}129 sessions). On the completed medium-scale run, dual LoRA gives the
best trained-model semantic Group@5/10 performance, while image-side LoRA
gives the strongest exact top-5 ranking. We argue that for a chat sticker
panel, semantic Group@5/10 is the right north-star metric and exact retrieval
should be treated as a secondary within-group ranking diagnostic.
\end{abstract}

%%%%%%%%% BODY TEXT
\section{Introduction}

When users reply with stickers in a chat, the choice depends on local
context (what was just said), long-range state (who is talking, prior tone,
recurring topics) and the communicative \emph{intent} behind the reply
(e.g.\ gratitude, encouragement, teasing). A retrieval system that only looks
at the last few turns frequently selects stickers with the wrong affect or the
wrong recurring style for a given user stream. Conversely, a system that uses
only image-side similarity ignores the dialogue entirely.

We focus on the U-Sticker corpus, which contains real multilingual chats
spanning crypto, gaming and social-domain channels. Several properties make
sticker retrieval different from generic image--text retrieval. First, the
candidate set is sparse and animated: 42{,}319 decodable stickers across
\texttt{.png}, \texttt{.gif} and \texttt{.webm}. Second, supervision is
\emph{single-positive}: each turn records the one sticker the user actually
sent, even though many other stickers in the bank are equally valid. Third,
turns frequently contain code-switching and chat slang outside CLIP's
pre-training distribution.

We build on the IGSR-style design~\cite{showandtell} of using a CLIP backbone
plus a small learned head, and extend it in three directions:
\begin{itemize}
\item A retrieved \emph{long-term memory} stream from previous sessions of the
      same user/file, encoded by multilingual~E5~\cite{wang2022e5}.
\item An LLM-generated \emph{reply intent} stream that abstracts away from the
      exact sticker id and describes the communicative function of the reply.
\item A unified pipeline that handles \texttt{.png}, \texttt{.gif} and
      \texttt{.webm} via per-frame OpenCLIP encoding and mean pooling, and
      supports parameter-efficient adaptation through LoRA~\cite{hu2021lora}
      on the text side, image side, or both.
\end{itemize}

\noindent\textbf{Contributions.}
(1)~A clean five-row experiment matrix (zero-shot CLIP, learned head with
memory, text-LoRA, image-LoRA, dual-LoRA) on a unified all-media sticker
bank.
(2)~A medium-scale benchmark showing that dual LoRA gives the most balanced
group-level retrieval, while image-side LoRA is strongest for exact top-5
ranking.
(3)~A reusable, self-contained MultiSticker codebase including data
manifest construction, memory retrieval, OpenCLIP wrapping, multi-frame
decoding, and SLURM-friendly run scripts.

%-------------------------------------------------------------------------
\section{Related Work}

\noindent\textbf{Vision--language retrieval.}
CLIP~\cite{radford2021clip} popularised dual-encoder image--text retrieval,
and OpenCLIP~\cite{ilharco2021openclip} provides open-weight reproductions
trained on LAION-class corpora. Our system uses ViT-B/32 with the
\texttt{laion2b\_s34b\_b79k} weights as the backbone for both the text and
image branches.

\noindent\textbf{Sticker / meme retrieval.}
Sticker reply suggestion has been studied as a multimodal retrieval and
ranking problem~\cite{laddha2020understanding,gao2020learning}. Intent-guided
sticker retrieval (IGSR)~\cite{zhang2022igsr} adds an explicit intent
classification head and a group prior. We adopt this template and add long-term
memory and parameter-efficient adaptation.

\noindent\textbf{Memory-augmented dialogue.}
Long-term dialogue memory has been used to condition open-domain
chatbots~\cite{xu2022long}. We use it specifically to construct a third query
stream for retrieval rather than to generate text. Our memory encoder is
multilingual~E5~\cite{wang2022e5}, chosen for code-switched chat data.

\noindent\textbf{Parameter-efficient fine-tuning.}
LoRA~\cite{hu2021lora} adds low-rank update matrices to selected linear
layers, leaving the original weights frozen. We inject LoRA into OpenCLIP's
\texttt{out\_proj}, \texttt{c\_fc} and \texttt{c\_proj} modules through
PEFT~\cite{mangrulkar2022peft}, and isolate text-side, image-side and dual
adaptation through name-based parameter masking.

%-------------------------------------------------------------------------
\section{Method}

\subsection{Pipeline overview}

\begin{figure*}[t]
\centering
\includegraphics[width=0.96\textwidth]{pipeline_overview.png}
\caption{End-to-end MultiSticker pipeline. Raw U-Sticker JSON is split into
12-hour sessions; an LLM produces a per-session memory and a per-sample
intent label/text; multilingual~E5 retrieves the top-3 historical sessions
for the current turn; OpenCLIP independently encodes the context, memory
and intent strings; the \texttt{IntentGuidedRetriever} head projects their
concatenation into the CLIP image space and ranks the sticker bank by
cosine similarity. An auxiliary intent classifier predicts the sticker's
group and supplies a prior for the two-stage rerank.}
\label{fig:pipeline}
\end{figure*}

For each turn $t$ where the user sent a sticker $s^\star_t$, we form three
text inputs: the immediate \emph{context} $c_t$ (up to 12 prior turns),
the retrieved \emph{memory} $m_t$ (concatenation of the top-3 prior session
summaries), and the LLM \emph{intent} $i_t$ (a short description of what the
sticker should express). These are encoded independently by the OpenCLIP
text encoder $f_T$. Sticker images $s$ are encoded by the OpenCLIP image
encoder $f_I$ (with mean-pooling over frames for animated formats). The
trainable head $g_\theta$ projects the concatenation into the image space:
\begin{equation}
\mathbf{q}_t = \mathrm{norm}\Bigl(g_\theta\bigl(
  [f_T(c_t);\,f_T(m_t);\,f_T(i_t)]
\bigr)\Bigr).
\end{equation}
Retrieval logits are
$\mathbf{z}_t = \mathbf{q}_t\,\mathbf{B}^{\top}/\tau$,
where $\mathbf{B}\!\in\!\mathbb{R}^{N\times d}$ is the (optionally adapted)
OpenCLIP image bank and $\tau$ is a learned temperature. A second head
predicts the sticker's intent group $y_t$, and we minimise
\begin{equation}
\mathcal{L} = \mathrm{CE}(\mathbf{z}_t, s^\star_t) +
              \lambda \,\mathrm{CE}(\mathbf{z}^{\mathrm{int}}_t, y_t),
\quad \lambda{=}0.2.
\end{equation}

\subsection{Long-term memory retrieval}

Each raw chat file is split into sessions by a 12-hour inactivity gap. For
sample $t$, we collect up to 8 prior sessions from the same file. Their
LLM-generated session summaries are encoded with the E5 \texttt{passage:}
prefix; the current session text is encoded with the \texttt{query:} prefix.
After L2 normalisation, the dot products rank the prior sessions and we keep
the top-3. The selected memory texts are concatenated and stored as $m_t$,
then encoded by OpenCLIP at training time. Crucially, memory enters the model
through the \emph{text} encoder, not as a separate vector, which lets the
retrieval head learn to reweight memory against the local context.

\subsection{LLM annotation}

We annotate every session with a $\le$~4-sentence
\texttt{session\_memory\_text} (used as the passage corpus for memory
retrieval), and every sticker-reply sample with one
\texttt{intent\_label} from a 10-class taxonomy
(celebration\_approval, humor\_teasing, encouragement\_support, gratitude,
affection\_care, confusion\_curiosity, frustration\_disapproval,
disappointment\_sadness, surprise\_reaction, neutral\_acknowledgment) plus a
free-form \texttt{intent\_text}. The intent label feeds the auxiliary head
and induces the semantic-group label space; the intent text becomes the
third query stream.

\subsection{Multi-frame sticker encoding}

\texttt{.png} stickers are encoded as one-frame clips. \texttt{.gif} and
\texttt{.webm} are decoded to all frames (FFmpeg for video, PIL for GIF),
each frame is encoded by the OpenCLIP image encoder, and the per-frame
embeddings are mean-pooled and renormalised. During image-LoRA training,
the full bank is re-encoded once per epoch without gradients; for each batch
we additionally re-encode only the unique \emph{positive} stickers with
gradients (animated positives capped at 4 frames) and patch them into a
cloned bank, which keeps backprop tractable for a 42K-sticker bank.

\subsection{Tuning modes}

We compare four tuning modes against a zero-shot CLIP baseline.
\textbf{direct\_clip} uses no learned head: the query is just $f_T(c_t)$
scored against the frozen image bank.
\textbf{head\_only} (\emph{memory}) trains the retrieval head on top of frozen
OpenCLIP encoders.
\textbf{text\_lora} additionally injects LoRA ($r{=}8$, $\alpha{=}16$) into
the text branch.
\textbf{image\_lora} injects LoRA into the visual branch only, while
\textbf{dual\_lora} adapts both branches. Only LoRA parameters and the head
are trainable; all original OpenCLIP weights are frozen.

%-------------------------------------------------------------------------
\section{Experiments}

\subsection{Setup}

We report the completed \texttt{rebuild\_medium} setting. The manifest is
constructed from 8{,}129 sessions and 368{,}503 provisional sticker-reply
samples. After frequency filtering and media decoding, the experiment uses
30{,}000 training samples, 5{,}000 validation samples (4{,}998 after media
filtering), and a 6{,}758-sticker bank over \texttt{.png}, \texttt{.gif}, and
\texttt{.webm}. All trainable models are optimized for 5 epochs with AdamW.
The retrieval head uses batch size 256; LoRA runs use batch size 96; inference
batch size is 256. The CLIP backbone is OpenCLIP ViT-B/32 with
\texttt{laion2b\_s34b\_b79k} weights. The retrieval temperature is $0.07$,
hidden dimension is 512, dropout is $0.1$, the auxiliary intent-loss weight
is $0.2$, and the intent vocabulary is clustered into 64 groups. LoRA uses
rank $r{=}8$, $\alpha{=}16$, dropout $0.05$, LoRA LR $10^{-4}$, and head LR
$10^{-3}$.

\noindent\textbf{Compared models.}
We evaluate five rows. \texttt{direct\_clip} is a zero-shot baseline with no
learned parameters. \texttt{head\_only} freezes OpenCLIP and trains only the
retrieval and intent heads. \texttt{text\_lora}, \texttt{image\_lora}, and
\texttt{dual\_lora} train the head plus LoRA adapters on the text branch, the
visual branch, or both branches respectively. This design isolates whether
domain adaptation is more useful on chat text, sticker imagery, or the shared
alignment geometry.

\subsection{Metrics}

We evaluate retrieval at two granularities. The primary metric is
\textbf{Group Hit@K}: a prediction is correct if any of the top-$K$ retrieved
stickers belongs to the same semantic intent group as the logged sticker. This
is the most faithful proxy for a chat sticker panel, where the interface shows
5--10 plausible choices and the user mainly needs the right affect, tone, or
communicative intent. We therefore treat \textbf{Group@5} and
\textbf{Group@10} as the north-star metrics. As a secondary diagnostic, we
also report \textbf{Exact Hit@K}, which requires recovering the single
observed sticker id. Exact metrics measure within-group ranking quality, but
are conservative because many unobserved stickers may be equally valid. All
Hit@K metrics are higher-is-better.

\subsection{Results}

\begin{table*}[t]
\setlength{\tabcolsep}{3.5pt}
\centering
\scriptsize
\begin{tabular}{l l c c c c c c c}
\toprule
\textbf{Model} & \textbf{Trainable}
 & \textbf{Group@1}$\uparrow$ & \textbf{Group@5}$\uparrow$
 & \textbf{Group@10}$\uparrow$ & \textbf{Group@30}$\uparrow$
 & \textbf{Exact@5}$\uparrow$ & \textbf{Exact@30}$\uparrow$ & \textbf{MAP}$\uparrow$ \\
\midrule
\texttt{direct\_clip} & none
 & 0.6335 & 0.8778 & \best{0.9084} & \best{0.9588} & 0.0012 & 0.0028 & 0.0013 \\
\texttt{head\_only} & head
 & \best{0.7933} & 0.8619 & 0.8780 & 0.9406 & 0.0296 & \second{0.1263} & 0.0282 \\
\texttt{text\_lora} & head+T
 & \second{0.7623} & 0.8661 & 0.8880 & 0.9420 & 0.0312 & 0.1190 & 0.0280 \\
\texttt{image\_lora} & head+V
 & 0.7063 & \second{0.8812} & 0.9044 & \second{0.9550} & \best{0.0376} & 0.1188 & \best{0.0302} \\
\texttt{dual\_lora} & head+T+V
 & 0.6275 & \best{0.8816} & \second{0.9064} & 0.9540 & \second{0.0326} & \best{0.1315} & \second{0.0286} \\
\bottomrule
\end{tabular}
\caption{\textbf{Main \texttt{rebuild\_medium} results.} The north-star
metrics are Group@5 and Group@10, matching a chat UI that displays a small
panel of sticker candidates. Exact metrics are stricter single-positive
diagnostics. \best{Red bold} marks the best method per column and
\second{blue bold} marks the runner-up. T/V denote text/visual LoRA.}
\label{table:medium-main}
\end{table*}

\begin{table*}[t]
\setlength{\tabcolsep}{5pt}
\centering
\scriptsize
\begin{tabular}{l l c c c c c}
\toprule
\textbf{Media} & \textbf{Model} & \textbf{Count}
 & \textbf{Group@5}$\uparrow$ & \textbf{Group@10}$\uparrow$
 & \textbf{Exact@5}$\uparrow$ & \textbf{Exact@30}$\uparrow$ \\
\midrule
PNG & \texttt{head\_only}  & 1{,}676 & 0.9141 & 0.9242 & 0.0197 & \best{0.1062} \\
PNG & \texttt{text\_lora}  & 1{,}676 & 0.9206 & 0.9314 & \second{0.0286} & \second{0.0990} \\
PNG & \texttt{image\_lora} & 1{,}676 & \second{0.9254} & \second{0.9391} & 0.0185 & 0.0752 \\
PNG & \texttt{dual\_lora}  & 1{,}676 & \best{0.9320} & \best{0.9403} & \best{0.0304} & 0.0973 \\
\midrule
GIF & \texttt{head\_only}  & 2{,}284 & 0.8695 & 0.8870 & \second{0.0499} & 0.1935 \\
GIF & \texttt{text\_lora}  & 2{,}284 & 0.8717 & 0.8971 & 0.0468 & 0.1843 \\
GIF & \texttt{image\_lora} & 2{,}284 & \best{0.8927} & \second{0.9168} & \best{0.0670} & \second{0.1948} \\
GIF & \texttt{dual\_lora}  & 2{,}284 & \second{0.8892} & \best{0.9177} & 0.0464 & \best{0.2062} \\
\midrule
WEBM & \texttt{head\_only}  & 1{,}038 & 0.7611 & 0.7832 & 0.0010 & 0.0106 \\
WEBM & \texttt{text\_lora}  & 1{,}038 & 0.7659 & 0.7977 & 0.0010 & 0.0077 \\
WEBM & \texttt{image\_lora} & 1{,}038 & \best{0.7842} & \second{0.8208} & \second{0.0039} & \best{0.0222} \\
WEBM & \texttt{dual\_lora}  & 1{,}038 & \second{0.7832} & \best{0.8266} & \best{0.0058} & \best{0.0222} \\
\bottomrule
\end{tabular}
\caption{\textbf{Per-media breakdown for trainable models.}
Direct CLIP is omitted because its result artifact does not store per-media
metrics. The group metrics are strongest and most stable in PNG/GIF, while
WEBM remains the hardest exact-retrieval format.}
\label{table:medium-media}
\end{table*}

Table~\ref{table:medium-main} and Table~\ref{table:medium-media} summarise
the completed medium-scale run. The key result is that \texttt{dual\_lora}
is the best trained model under the panel-recommendation objective: it wins
overall Group@5 and is second only to the zero-shot baseline on Group@10.
At the same time, \texttt{image\_lora} has the strongest exact top-5 ranking,
which suggests that visual adaptation is especially useful for ordering
near-duplicate stickers within a semantic group.

%-------------------------------------------------------------------------
\section{Analysis and Discussion}

\noindent\textbf{Group retrieval is the right north star.}
The most important deployment regime is not auto-sending a single sticker;
it is filling a chat panel with 5--10 plausible choices. In this regime,
recovering the \emph{semantic group} of the desired reply is more important
than recovering the exact historical sticker id. The results support this
framing: exact retrieval remains difficult for all methods, but Group@5 and
Group@10 are already near the level required for a usable recommendation
panel. Under this criterion, \texttt{dual\_lora} is the best trained model:
it achieves the strongest Group@5 (0.8816) and the strongest trained-model
Group@10 (0.9064). The zero-shot \texttt{direct\_clip} baseline slightly
exceeds it on Group@10, but its exact ranking is nearly unusable, indicating
that CLIP recognises coarse semantic neighborhoods while failing to
distinguish the concrete sticker vocabulary.

\noindent\textbf{Exact metrics are secondary but still diagnostic.}
Exact Hit@K should not be treated as the primary objective because U-Sticker
contains only one logged positive per context. A different sticker with the
same intent may be a perfectly good recommendation but is counted as wrong
by Exact Hit@K. Still, exact metrics reveal how well a model orders siblings
inside a group. Here \texttt{image\_lora} has the best Exact@5 and MAP,
suggesting that visual adaptation improves fine-grained ordering among
similar stickers. \texttt{dual\_lora}, meanwhile, has the best Exact@30,
which is attractive for a two-stage system: first retrieve a semantically
correct candidate pool, then rerank or diversify it for display.

\noindent\textbf{Why LoRA helps different media differently.}
The per-media table shows that GIF is the most responsive format to visual
adaptation: \texttt{image\_lora} wins GIF Exact@5 and Group@5, while
\texttt{dual\_lora} wins GIF Exact@30 and Group@10. This is intuitive:
animated stickers encode much of their intent in character pose, motion, and
style, so adapting the visual branch makes the embedding space more
sticker-specific. PNG behaves differently. The frozen CLIP space already
separates many static stickers well, so \texttt{head\_only} remains strong
on PNG Exact@30, while \texttt{dual\_lora} gives the best PNG Group@5/10.
WEBM remains the hardest exact-retrieval case; frame sampling and video
compression likely erase cues needed to distinguish close variants.

\noindent\textbf{Recommended retrieval design.}
For a production-style sticker panel, the cleanest design is to use
\texttt{dual\_lora} as the group-aware first-stage retriever and optimize
the system around Group@5/10. Exact@5 and MAP should be tracked as secondary
ranking signals, where \texttt{image\_lora} currently sets the strongest
reference point. A future reranker can combine these strengths: use
\texttt{dual\_lora} to guarantee semantic coverage, then apply a lightweight
visual/textual reranker or diversity rule to choose the final 5--10 stickers.

%-------------------------------------------------------------------------
\section{Conclusion}

We presented MultiSticker, an intent-guided dialogue-aware sticker
retrieval system that combines OpenCLIP, retrieved long-term memory, and
LLM-generated intent into a single trainable head with optional LoRA
adaptation on either encoder branch. On the completed medium-scale U-Sticker
run, dual LoRA is the strongest trained model for group-level panel
recommendation, while image-side LoRA gives the best exact top-5 ordering.
This supports a two-stage view of sticker retrieval: optimize the first stage
for semantically correct 5--10 item panels, then use exact metrics as
secondary diagnostics for fine-grained reranking. The released code makes the
five-row comparison reproducible end-to-end.

\noindent
{\bf Acknowledgments.} We thank the authors of OpenCLIP and PEFT, and the
U-Sticker maintainers. Session summaries and intent annotations were
produced with Qwen2.5-32B-GPTQ; all experimental decisions are our own.

{\small
\bibliographystyle{ieee}
\bibliography{egbib}
}

\end{document}
